% FONTE TEMA https://github.com/matze/mtheme
\documentclass[aspectratio=1610]{beamer}
%\documentclass[aspectratio=1610, handout]{beamer}
\usepackage[utf8]{inputenc}
\usepackage{ragged2e}
\usepackage{xcolor}
\usepackage[italian]{babel}
\usepackage{multirow}
\usepackage{silence}
\usepackage{listings}
\definecolor{verde}{rgb}{0,0.6,0}
\definecolor{grigio}{rgb}{0.5,0.5,0.5}
\definecolor{rosa}{rgb}{0.58,0,0.82}
\definecolor{backgroud}{rgb}{0.95,0.95,0.92}
\lstdefinestyle{stilepython}{
    backgroundcolor=\color{backgroud},
    commentstyle=\color{verde},
    keywordstyle=\color{magenta},
    numberstyle=\tiny\color{grigio},
    stringstyle=\color{rosa},
    basicstyle=\ttfamily\footnotesize,
    breakatwhitespace=false,
    numbers=left,
    numbersep=5pt,
    showspaces=false,
    showstringspaces=false,
    showtabs=false,
    literate=
    {à}{{\'a}}1
    {è}{{\'e}}1
    {ì}{{\'i}}1
    {ò}{{\'o}}1
    {ù}{{\'u}}1
}
\lstset{style=stilepython}

\WarningFilter{beamer}{}
\WarningFilter{metropolis}{}
\usetheme[progressbar=frametitle,titleformat=smallcaps]{metropolis}
\setbeamertemplate{frame numbering}[fraction]
\setbeamercovered{dynamic}
\definecolor{rosso}{RGB}{255, 0, 0}
\definecolor{giallo}{RGB}{254,212,23}
\hypersetup{colorlinks=true,linkcolor=black,urlcolor=rosso}
\setbeamercolor{palette primary}{fg=black, bg=giallo}
\setbeamercolor{background canvas}{bg=white}
\setbeamercolor{normal text}{fg=black}
\setbeamercolor{progress bar}{fg=rosso}
\setbeamercolor{framesubtitle}{fg=rosso}
\setbeamercolor{normal text .dimmed}{fg=giallo}
\setbeamercolor{block title alerted}{fg=rosso, bg=giallo}
\setbeamerfont{caption}{size=\tiny}
\setbeamerfont{caption name}{size=\tiny}
\setlength{\abovecaptionskip}{0pt}
\makeatletter
\metroset{block=fill}
\setlength{\metropolis@progressinheadfoot@linewidth}{1pt} 
\setlength{\metropolis@progressonsectionpage@linewidth}{1pt}
\setlength{\metropolis@titleseparator@linewidth}{1pt}
\makeatother

\title{SOVRANIT\'A DIGITALE}
\subtitle{Dominio tecnologico vs indipendenza tecnologica}
\date{}
\institute{\textit{
        Fonti:
        \begin{itemize}
            \item[-] \href{www.nomefonte1.it}{NOME FONTE 1}
            \item[-] \href{www.nomefonte2.it}{NOME FONTE 2}
        \end{itemize}
    }
}

\begin{document}

\begin{frame}[plain, noframenumbering]
    \titlepage
\end{frame}

\begin{frame}{SOVRANIT\'A DIGITALE VS CAPITALISMO DELLE PIATTAFORME}
    \begin{alertblock}{SOVRANIT\'A DIGITALE}
        \begin{minipage}{0.98\linewidth}
            \justifying
            La \textbf{sovranità digitale} è la capacità di uno Stato e dei suoi cittadini di 
            \textbf{mantenere il controllo autonomo sui propri dati e sulle infrastrutture informatiche}.\\
            \bigskip
        \end{minipage}
    \end{alertblock}
    \pause
    \begin{alertblock}{CAPITALISMO DELLE PIATTAFORME}
        \begin{minipage}{0.98\linewidth}
            \justifying
            Poche \textbf{multinazionali} (Amazon, Apple, ByteDance, Google, Meta, Microsoft, SpaceX) extra-europee 
            \textbf{controllano l'infrastruttura fondamentale} su cui poggia l'economia e la comunicazione.\\
            \bigskip
        \end{minipage}
    \end{alertblock}
\end{frame}

\end{document}